\subsection{退化面の元残余を積として使う}
\label{ndg25:section}
$a=uh-B(F+1)=0$ において，$A=Bs-b$，$R=b^2-sT$，
$L=(F+1)R-hA$，$M=uR-BA$ と置くと
\[
 hM=BL,\quad FL=bF+s(T-h),\quad
 \mathfrak r=M(t-y)-tT,\quad M>0.
\]
元 pair 方程式により $M(t-y)=tT$ であり，退化面上では $y<t$，$M\ge T$，$M\mid tT$ が必要になる．
一般の三高模型では $j=t-y$ の正値性を仮定してはいけない．

$h\ge b$ なら $T<h$，$1\le L\le b-1$，$h\mid(b-1)L$ なので
$\gcd(h,b-1)\ge2$ である．従って互いに素な領域を上限なしに除ける．
$x\ge1$ の退化面では $s-bx>0$ が $h\ge b$ を保証する．
さらに
\[
 b\ge3,\quad b+1\le s\le b^2-1,\quad
 1\le x\le\lfloor(s-1)/b\rfloor,\quad 1\le T\le M\le b-2.
\]
固定 $b$ から全残りの整数を有限復元できるので，従来の固定 $s$ からの復元と相補的に使う．
先に $g=\gcd(h,b-1)$ を計算し，$L$ を $h/g$ の倍数だけに絞ると無駄が少ない．
積表示は残った非互いに素な場合の整除性や共通因子を調べる入口になる．
この帰着は $b$ の一様上界や全退化面の排除を意味しない．
